\section{Data Access and Retrieval}
\label{sec:data-access}

\subsection{Source Files}

The IHPD is published by the Globalization Institute of the Federal Reserve
Bank of Dallas at
\begin{center}
  \url{https://www.dallasfed.org/research/international/houseprice}
\end{center}

Each release consists of two Excel workbooks, described in
Table~\ref{tab:release-files}. The filename encodes the last quarter
included, not the release date: \texttt{hp2504.xlsx} is the vintage whose
final observation is 2025\,Q4, published in the first full week of January
2026 under the schedule in Table~\ref{tab:release-schedule}.

\begin{table}[ht]
  \centering
  \caption{IHPD quarterly release files}
  \label{tab:release-files}
  \begin{tabular}{lp{9cm}}
    \toprule
    File & Contents \\
    \midrule
    \texttt{hp\textit{YYQ}.xlsx}
      & The data. One worksheet per variable --- \texttt{HPI},
        \texttt{RHPI}, \texttt{PDI}, \texttt{RPDI} --- with quarters in rows
        and economies in columns, plus the two aggregate columns. A
        \texttt{Note} sheet carries citation details and flags which series
        contain forecast values. \\[4pt]
    \texttt{hpta\textit{YYQ}.xlsx}
      & Pre-computed exuberance statistics: GSADF and SADF statistics with
        95\,\% critical values by country, and BSADF sequences for two
        indicators (real house prices, and the ratio of RHPI to RPDI).
        Separate sheets report results for lag orders 1 and 4. \\
    \bottomrule
  \end{tabular}
\end{table}

\noindent Note that the layout is by variable, not by country: a single
worksheet holds all economies for one variable. The BSADF sequences in the
\texttt{hpta} file begin later than the data, because the recursive procedure
requires a minimum initial estimation window (Section~\ref{subsec:gsadf}).

Only the current release is linked from the Dallas Fed page. Earlier
vintages are retained by the Dallas Fed, and the Observatory additionally
archives the release files it has used.

\subsection{The Observatory Data API}
\label{subsec:api}

The Observatory publishes its own processed outputs --- not only the source
data, but the estimated statistics behind the dashboard --- as JSON over
HTTP, versioned by release:
\begin{center}
  \url{https://api.housing-observatory.com/datasets/int/index.json}
\end{center}

The index document lists available releases, most recent first. Each release
then exposes the series and estimation objects used by the dashboard, for
example \texttt{rhpi.json} and \texttt{pti.json} for the two indicator
series, and the corresponding fitted test objects. This is the interface the
dashboard itself consumes, so anything visible on the platform can be
reproduced from it without re-running the estimation.

\subsection{Quarterly Update Workflow}

Each release is incorporated through a fixed sequence:

\begin{enumerate}
  \item The Dallas Fed publishes the new \texttt{hp} and \texttt{hpta}
        workbooks on the schedule in Table~\ref{tab:release-schedule}.
  \item The release is retrieved from the Dallas Fed and archived.
  \item The two indicator series are constructed: the real house price index,
        and the price-to-income ratio RHPI/RPDI.
  \item Exuberance statistics are re-estimated with \texttt{exuber}
        (Section~\ref{subsec:gsadf}) and the PSY-IVX decomposition is
        re-estimated for the countries it covers
        (Section~\ref{subsec:psyivx}).
  \item Results are written to the data API
        (Section~\ref{subsec:api}), from which the dashboard reads.
  \item The quarterly monitoring report and the nowcast report are
        regenerated and published.
  \item This report is updated and its changelog
        (Appendix~\ref{app:changelog}) extended.
\end{enumerate}
